\documentclass[11pt,a4paper]{article}
\usepackage[T2A,T1]{fontenc}
\usepackage[utf8]{inputenc}
\usepackage[english,russian]{babel}
\usepackage{geometry}
\usepackage{graphicx}
\usepackage{booktabs}
\usepackage{microtype}
\usepackage[hidelinks]{hyperref}
\geometry{margin=20mm}
\graphicspath{{results_tau1/}{results/}}
\setlength{\parindent}{1.1em}
\setlength{\parskip}{0.2em}
\newcommand{\code}[1]{{\fontencoding{T1}\selectfont\ttfamily #1}}

\title{EDM-анализ Muon/nanoGPT: фиксированная задержка $\tau=1$ и четыре метода}
\author{Технический отчёт по \code{message (2).txt}}
\date{30 июля 2026 г.}

\begin{document}
\maketitle

\begin{abstract}
Для плотного ряда training loss одного запуска Muon/nanoGPT проведён
скользящий анализ с фиксированной задержкой $\tau=1$. Использованы все четыре
метода размерностного анализа из существующего проекта: False Nearest
Neighbors, Cao, simplex projection и Levina--Bickel MLE. Методы не дают
согласованного свидетельства позднего dimensionality collapse: FNN слабо
снижается, в то время как Cao, simplex и MLE в среднем растут или остаются
высокими. Результат является negative/control case, а не свидетельством
grokking-перехода.
\end{abstract}

\section{Данные и режим расчёта}

Из \code{icomp\_article/message (2).txt} извлечён один завершённый запуск
\code{Muon\_lr0.06\_5db64adc}: 2330 последовательных значений training loss и
11 validation-loss точек. Training loss в исходном логе является усреднённой
между GPU суммой loss локального батча; для анализа он нормализован к token-level
cross-entropy делением на 32768. Validation loss используется лишь как внешний
индикатор: 11 точек недостаточно для корректной фазовой реконструкции.

Во всех окнах длиной 400 шагов со сдвигом 40 использована фиксированная
задержка $\tau=1$ и максимальная размерность embedding $E_{max}=15$. Методы
вызываются непосредственно из \code{code/Grokking/search\_for\_optimal\_parameters.py}
с настройками существующего pipeline: FNN выбирает первое $E$ с долей false
neighbours ниже 1\%, Cao использует первое $|\Delta E_1|<0.05$, simplex выбирает
минимальный one-step out-of-sample RMSE, а MLE использует $k=5$ соседей.

Пунктирная вертикаль на рисунках --- начало learning-rate cooldown (шаг около
1260), точечная --- начало 40-шагового extension (2290).

\begin{figure}[ht]
\centering
\includegraphics[width=\textwidth]{01_training_overview.pdf}
\caption{Training и validation loss исходного запуска вместе с learning-rate schedule.}
\end{figure}

\section{Четыре метода при $\tau=1$}

\begin{figure}[ht]
\centering
\includegraphics[width=\textwidth]{tau1_all_methods_raw.pdf}
\caption{Скользящие оценки четырьмя legacy-методами на исходном train loss.}
\end{figure}

FNN даёт значения около 3--7 и слабое позднее снижение. Cao имеет постоянное
$E=2$ до cooldown, затем преимущественно принимает значения 6--8 и возвращается
к 2 в последних окнах. Simplex имеет очень неустойчивый аргумент минимума RMSE
и часто выбирает верхнюю границу $E=15$. MLE непрерывно растёт с ранних значений
около 11--13 к поздним 14--16.

\begin{table}[ht]
\centering
\caption{Средние первых и последних 12 скользящих окон. Изменение равно late минус early.}
\small
\begin{tabular}{lrrrrrr}
\toprule
& \multicolumn{3}{c}{Raw loss} & \multicolumn{3}{c}{Detrended loss} \\
Метод & Early & Late & $\Delta$ & Early & Late & $\Delta$ \\
\midrule
FNN & 4.42 & 4.17 & -0.25 & 4.17 & 3.67 & -0.50 \\
Cao & 2.00 & 4.25 & +2.25 & 5.58 & 6.33 & +0.75 \\
Simplex & 10.08 & 13.17 & +3.08 & 10.00 & 12.00 & +2.00 \\
MLE ID & 12.88 & 14.77 & +1.89 & 14.39 & 15.30 & +0.91 \\
\bottomrule
\end{tabular}
\end{table}

\begin{figure}[ht]
\centering
\includegraphics[width=\textwidth]{tau1_raw_vs_detrended.pdf}
\caption{Сравнение исходного и локально detrended loss. Удаление тренда не
создаёт согласованного падения оценок.}
\end{figure}

\section{Критерии выбора размерности}

\begin{figure}[ht]
\centering
\includegraphics[width=\textwidth]{tau1_method_criteria.pdf}
\caption{Критериальные кривые FNN, Cao и simplex в representative early, middle
и late окнах. MLE не выбирает дискретное $E$: он оценивает локальную ID в
пространстве размерности 15.}
\end{figure}

FNN --- пороговая дискретная эвристика, поэтому снижение на 0.25--0.50 не
является сильным структурным эффектом. В поздних окнах процент false neighbours
не достигает устойчивого нуля при малых $E$. Cao чувствителен к своему
plateau-правилу, что видно по скачку от 2 к 6--8 в cooldown. У simplex минимум
RMSE часто лежит на границе допустимого диапазона, следовательно, его значения
нельзя читать как устойчивую оценку истинной размерности. MLE не демонстрирует
падения ни для raw, ни для detrended ряда.

\section{Вывод}

Четыре метода отвечают на разные операциональные вопросы, поэтому их
абсолютные значения сравнивать напрямую нельзя. Однако для гипотезы о
dimensionality collapse важно согласованное направление, а его здесь нет: один
пороговый метод слегка снижается, три других не подтверждают этот сигнал.

Следовательно, этот одиночный Muon/nanoGPT run следует использовать как
negative/control case: непрерывное успешное обучение LLM не обязано порождать
универсальный поздний коллапс оценок размерности, построенных только по
одномерному training loss. Для сильного сравнительного вывода необходимы
повторы по seed и логи остальных оптимизаторов того же benchmark.

\end{document}
